\section{Conclusion}

This paper reframes UGDA evaluation around a deployment requirement: selecting
an algorithm, configuration, seed, and checkpoint without reading target
labels. \benchmark makes this requirement auditable through immutable
trajectories and sealed labels, while \method estimates target risk with a
covariance-aware recovery system built on candidate disagreement. The tight
graph-spectral frame supplies an exact decomposition and exposes band-wise
heterogeneity, but the resulting analysis identifies cross-model error
covariance as the central hidden term and quantifies how covariance transport
error affects recovered risk. In the
current frozen development scope, the method lowers mean normalized regret by
$37.9\%$, ranks the candidate bank accurately near its optimum, and satisfies
the zero-access protocol on source-simulated folds. On the separately sealed
real-target Gate-1 comparison, however, Transfer Score remains stronger in
top-1 regret ($0.1467$ versus $0.2560$), so the current evidence does not
support a state-of-the-art claim.

The core ablation sharpens the technical conclusion. Covariance correction is
essential on source-simulated folds, while spectral band conditioning provides
only a modest, statistically inconclusive improvement over an equal-information
global covariance control. Robust refitting and bootstrap uncertainty do not
drive the result and are therefore implementation details, not contributions.

The central limitation is transportability: current unlabeled descriptors do
not reliably predict risk-correction geometry for unseen shift families. The
next method iteration should therefore be conservative rather than more
expressive: penalize transport uncertainty, fuse the recovered risk with an
independent transferability prior such as Transfer Score, and choose all
weights by leave-one-shift-family-out validation. This reliability question
must be addressed before the frozen 16-transfer evaluation, without additional
target-specific tuning.
